\documentclass[hyperref={pdfpagelabels=false},table]{beamer}

\input{lec_common}
\begin{document}

\part{Unit 6: Object oriented programming with classes}

%\frame{\tableofcontents}





\begin{frame}[fragile]{Datatypes and Classes}
A datatype, e.g. a \pyth!list! binds data and related methods
together:

\begin{python}
my_list= [1,2,3,4,-1]
my_list.<tab> #  show a list of all related methods (in ipython only)
my_list.sort()  # is such a method
my_list.reverse() # .. another method
\end{python}

In this unit we show how own dataypes and related methods can be created.

For example
\begin{itemize}
\item polynomials
\item triangles
\item special problems (find a zero)
\item etc.
\end{itemize}
\end{frame}
\begin{frame}[fragile]{A minimalistic example}
\begin{python}
class Nix:
    pass
\end{python}

An object  with a certain datatype is said to be
an \emph{instance} of that type:

\begin{python}
a=Nix()
\end{python}

Check:
\begin{python}
 if isinstance(a,Nix):
    print('Indeed it belongs to the new class Nix')
\end{python}
\end{frame}

%%%%%%%%%%%%%%%%%%%%%%%%%%%%%%%%%%%%%%%%%%%%%%%%%%%%%%%%%%%%%%%%%%%%%
%%%%%%%%%%%%%%%%%%%%%%%%%%%%%%%%%%%%%%%%%%%%%%%%%%%%%%%%%%%%%%%%%%%%%
\begin{frame}[fragile]\frametitle{Naming Convention}

  See \href{https://www.python.org/dev/peps/pep-0008/}{PEP 8: Style Guide for Python Code} for details.

  Function and variable names should be separated with underscore.
  \begin{python}
    my_special_variable = 1 # variables 

    def my_special_function():
      pass
  \end{python}

  \pause 
  Or use \pyth{CapitalizedWords} with initial lower case character. 
  \begin{python}
    mySpecialVariable = 1    # variables 

    def mySpecialFunction(): # functions
      pass
  \end{python}

  and classes like this 

  \begin{python}
    class MySpecialClass:    # classes
        ...
  \end{python}

\end{frame}


\begin{frame}[fragile]{A more complete example}
Let us define a new datatype for rational numbers:
\begin{python}
class RationalNumber:
    def __init__(self,numerator,denominator):
        if not isinstance(numerator,int) :
             raise TypeError(
                   "numerator should be of type int")
        if not isinstance(denominator,int) :
             raise TypeError(
                   "denominator should be of type int")
        self.numerator=numerator
        self.denominator=denominator
\end{python}
\vfill
It got a method which initializes an instance with two attributes:
\vfill
\begin{python}
q=RationalNumber(10,20)    # Defines a new object
q.numerator                # returns 10
q.denominator              # returns 20
p=RationalNumber(10,20.0)  # raises a TypeError
\end{python}

\end{frame}
\begin{frame}[fragile]{\pyth!__init__! and \pyth!self!}
What happens, when
\begin{python}
q=RationalNumber(10,20)    # Defines a new object
\end{python}
is executed?
\vfill
\begin{itemize}
\item a new object with name \pyth!q! is created
\item the command \pyth!q.__init__(10,20)! is executed.
\end{itemize}
\vfill
\pyth!self! is a placeholder for the name of the newly created instance (here: \pyth!q!)
\vfill

\end{frame}
\begin{frame}[fragile]{Adding methods}
Methods are functions bound to the class:
\vfill
\begin{python}
class RationalNumber:
...
  def convert2float(self):
    return float(self.numerator)/float(self.denominator)
\end{python}
\vfill
and its use ...
\begin{python}
q=RationalNumber(10,20)   # Defines a new object
q.convert2float()         # returns 0.5
\end{python}
\vfill
Note, again the special role of \pyth!self! !
\end{frame}
\begin{frame}[fragile]{Adding methods\hfill (Cont.)}
\vfill
Note: \\
Both commands are equivalent!
\begin{python}
RationalNumber.convert2float(q)
q.convert2float()
\end{python}
\vfill
\end{frame}
\begin{frame}[fragile]{Adding method\hfill(Cont.)}
A second method ....
\begin{python}
class RationalNumber:
...
  def add(self, other):
    p1, q1 = self.numerator, self.denominator
    if isinstance(other, RationalNumber):
      p2, q2 =  other.numerator, other.denominator
    elif isinstance(other, int):
      p2, q2 = other, 1
    else:
      raise TypeError("...")
    return RationalNumber(p1 * q2 + p2 * q1, q1 * q2)
\end{python}

A call to this method takes the following form

\begin{python}
q = RationalNumber(1, 2)
p = RationalNumber(1, 3)
q.add(p)   # returns the RationalNumber for 5/6
\end{python}

\end{frame}
\begin{frame}[fragile]{Special methods: Operators}
We would like to add RationalNumber number instances
just by\\
\centerline{ \pyth!q+p!.}

Renaming the method\\
\vfill
\begin{python}
RationalNumber.add
\end{python}
 \centerline{to} 
\begin{python}
RationalNumber.__add__
\end{python}
\vfill
makes this possible.
\end{frame}
\begin{frame}[fragile]{Special methods: Operators \hfill(Cont.)}
\begin{tabular}{ll|ll}
\hline
Operator & Method&Operator&Method\\
\hline
{\pyth!+!} & {\pyth!__add__!}&{\pyth!+=!}&{\pyth!__iadd__!}\\
{\pyth!*!} & {\pyth!__mul__!}&{\pyth!*=!}&{\pyth!__imul__!}\\
{\pyth!-!} & {\pyth!__sub__!}&{\pyth!-=!}&{\pyth!__isub__!}\\
{\pyth!/!} & {\pyth!__truediv__!}&{\pyth!/=!}&{\pyth!__idiv__!}\\
{\pyth!**!} & {\pyth!__pow__!}&&\\
\hline
{\pyth!==!} & {\pyth!__eq__!}&{\pyth|!=|}&{\pyth!__ne__!}\\
{\pyth!<=!} & {\pyth!__le__!}&{\pyth|<|}&{\pyth!__lt__!}\\
{\pyth!>=!} & {\pyth!__ge__!}&{\pyth|>|}&{\pyth!__gt__!}\\
\hline
{\pyth!()!} & {\pyth!__call__!}&{\pyth|[]|}&{\pyth!__getitem__!}\\
\hline
\end{tabular}
\end{frame}
\begin{frame}[fragile]{Special methods: Representation}
\pyth!__repr__! defines how the object is represented,
when just typing its name:

\begin{python}
class RationalNumber:
 ...
  def __repr__(self):
    return f"{self.numerator} / {self.denominator}"
\end{python}
\vfill
Now:
 \pyth!q! returns \pyth!10/20!.

\end{frame}
\begin{frame}[fragile]{Reverse operations}
So far we can use the class for operations like
\[
1/5+5/6  \texttt{ or } 1/5 + 5
\]
but $5 + 1/5$ requires that the integer's method \pyth!__add__!
knows about \pyth!RationaNumber!.
\vfill
Instead of extending the methods of \pyth!int! we use reverse operations:
\begin{python}
class RationalNumber(object)
   ....
    def __radd__(self, other):
        return self + other
\end{python}
\vfill
\pyth!__radd__! reverses the role of self and other.
\vfill
Check the operation $1 + q$
\end{frame}
\begin{frame}[fragile]{Extending the example}
We add a method \pyth!shorten! to our example:
\vfill
\begin{python}

class RationalNumber:
    ....
    def shorten(self):
        def gcd(a, b):
            # Computes the greatest common divisor
            if b == 0:
                return a
            else:
                return gcd(b, a % b)
        factor = gcd(self.numerator, self.denominator)
        self.numerator = self.numerator // factor
        self.denominator = self.denominator // factor
\end{python}
\vfill
Note, it performs an \emph{inplace} modification of the object.
\vfill
\end{frame}


\end{document}

